% TOP_PROBLEM: {"id": 3178, "title": "Hamilton decomposition of prisms over 3-connected cubic planar graphs", "category_id": 3, "category": {"id": 3, "name": "graph_theory", "display_name": "Graph Theory", "description": "Problems involving graphs, networks, and their properties.", "slug": "graph-theory", "order_index": 3, "created_at": "2026-07-31T15:26:25.671Z"}, "status": "open", "problem_number": "OPG-47356", "set_id": 12, "statusQualification": "The complementary cycle factor is taken in the prism G square K2; the source discussion incorrectly writes the base graph G minus E(C)."}
% FIRST_POSTED: 2026-10-01
% ENTRY_KIND: statement_only
% UnsolvedMath ID 3178; source catalogue code OPG-47356
% !TeX program = lualatex
% Definition and sources only; this file is not an LLM solution attempt.
\documentclass[11pt,a4paper]{article}
\usepackage[margin=25mm]{geometry}
\usepackage{fontspec}
\setmainfont{lmroman10-regular.otf}[BoldFont=lmroman10-bold.otf,
 ItalicFont=lmroman10-italic.otf,BoldItalicFont=lmroman10-bolditalic.otf]
\usepackage{amsmath,amssymb,mathtools}
\usepackage{unicode-math}
\setmathfont{latinmodern-math.otf}
\AtBeginDocument{\let\setminus\smallsetminus}
\usepackage{xurl,hyperref}
\hypersetup{colorlinks=true,urlcolor=blue}
\setcounter{secnumdepth}{0}
\setlength{\parindent}{0pt}
\setlength{\parskip}{5pt}
\setlength{\emergencystretch}{3em}
\begin{document}
\section*{Hamilton decomposition of prisms over 3-connected cubic planar graphs}\label{3178}
{\small UnsolvedMath ID 3178 · OPG-47356 · Graph Theory · OpenGarden}
\subsection{Definitions and mathematical statement}
Let $G=(V,E)$ be a finite simple $3$-connected cubic planar
graph. Thus every vertex has degree three, $G$ has a
crossing-free drawing in the plane, and deletion of any
two vertices leaves it connected. Define its prism
$P=G\square K_2$ on vertex set $V\times\{0,1\}$.
For each $uv\in E$ and $i\in\{0,1\}$ put an edge
$(u,i)(v,i)$ in $P$, and for each $v\in V$ put the
vertical edge $(v,0)(v,1)$ in $P$.

The prism is a $4$-regular graph with $2|V|$ vertices.
A Hamilton decomposition into two cycles means Hamilton
cycles $C_1,C_2$ in $P$ such that
\[
 E(P)=E(C_1)\mathbin{\dot\cup}E(C_2).
\]
Each cycle must visit all vertices of both layers, and
each prism edge must be used exactly once by the pair.

The conjecture is that every base graph $G$ satisfying
the stated hypotheses has such a decomposition. Equivalently,
one seeks a Hamilton cycle $C$ of $P$ whose edge complement
$(V(P),E(P)\setminus E(C))$ is connected. That complement
automatically has degree two at every vertex; connectivity
is what makes it one Hamilton cycle instead of several
disjoint cycles. The complement is taken in the prism $P$,
not in the base graph $G$. A Hamilton cycle in each layer
separately would not satisfy the spanning requirement.
\subsection{Short English statement}
Can the prism over every 3-connected cubic planar graph be partitioned into two Hamilton cycles?
\subsection{Sources}
\begin{itemize}
\item[S1] ULAM.AI and the UnsolvedMath Contributors, supplied problems.json, record 3178 (OPG-47356), OpenGarden collection. Catalogue identity and source transcription; CC BY 4.0.
\url{https://huggingface.co/datasets/ulamai/UnsolvedMath}.
\item[S2] Open Problem Garden, Hamilton decomposition of prisms over 3-connected cubic planar graphs. Original problem and discussion; the mathematical formulation below is expanded from the supplied source record.
\url{http://www.openproblemgarden.org/op/decomposing_the_prism_of_a_3_connected_cubic_planar_graphs_in_hamilton_cycles}.
\end{itemize}
\end{document}
